% docs/manuscript/sections/prop3_v2.tex
% Stream S2, deliverable 1. Appendix: the domain of Proposition 3, W as a random variable,
% and the proofs of thm:floor (tightened), prop:invariance and cor:split.
% Every numeral is read from results/S2_theory/tables/ or results/S6_synchronized_traces/tables/.
% Produced by experiments/S2_theory/{s2_arl0.py, s2_w_random.py, s2_eddm.py}.

\section{Proposition 3: Domain, Randomness of \texorpdfstring{$W$}{W}, and Proofs}\label{sec:prop3}

\subsection{Where Eq.~\eqref{eq:markov_bound} says anything}\label{sec:prop3_domain}

Proposition~\ref{prop:starvation} is sound: Azuma--Hoeffding applied to the
centred walk, Doob's maximal inequality on the exponential submartingale, a
union bound over the at most $W$ reflection restart points, with the
$\sqrt{W \ln(W/\varepsilon)}$ fluctuation margin retained and $s_0$ carried
explicitly. What was never stated is the set on which it is not vacuous.

\begin{definition}[Informative domain]\label{def:eq4_domain}
  Eq.~\eqref{eq:markov_bound} bounds a probability, so its content is
  $\min\{1,\, W e^{-\frac{2}{W}(\lambda - s_0 - \mu W)_+^2}\}$. Write
  \[
    \mathcal{D} \;:=\; \bigl\{ (W, \lambda) \;:\; \mu W < \lambda - s_0 \bigr\},
    \qquad \mu := \Delta e - \delta_P .
  \]
  Outside $\mathcal{D}$ the positive part vanishes, the exponential is $1$, and
  the bound returns $W \ge 1$: it is \emph{vacant}. Membership of $\mathcal{D}$
  is necessary but not sufficient --- $W e^{-\frac{2}{W}(\cdot)^2}$ must also
  fall below $1$ --- and the two criteria are reported separately below.
\end{definition}

Instantiated at $\delta_P = \DeltaPtext$, $s_0 = 0$, on the three measured
estimators of $W$ (mean over $100$ seeds per magnitude, arm \texttt{full}) and
the threshold ladder $\lambda \in \{8, 15, 25, 50\}$, the $240$ pairs split
$188$ vacant and $52$ inside $\mathcal{D}$, of which $32$ yield a bound
strictly below $1$. At the canonical magnitude $\Delta e = 0.3268$:

\smallskip
\centerline{\begin{tabular}{lrrrrr}
  \toprule
  $W$ estimator & mean $W$ & $\mu W$ & $\lambda = 15$ & $\lambda = 25$ & $\lambda = 50$ \\
  \midrule
  $\tau_{\mathrm{swap}}^{(1/M)}$ & 57.4  & 18.2  & vacant & $1.0$ & $2.7\!\times\!10^{-14}$ \\
  $\tau_{\mathrm{erase}}$        & 611.9 & 193.8 & vacant & vacant & vacant \\
  $\tau_{\mathrm{err}}(0.10)$    & 557.0 & 176.4 & vacant & vacant & vacant \\
  \bottomrule
\end{tabular}}
\smallskip

\noindent This is Remark~\ref{rem:transient_length} verified rather than
asserted: $\mu W = \SixMuWfirst$ at the first replacement and $\SixMuWerase$ at
erasure reproduce from the campaign to the printed precision. The remark's
diagnosis is also correct and is the one that matters --- the hypothesis that
fails is the \emph{constant accumulation rate} $\mu$ across the window, not the
fluctuation term. The error excess decays inside the window; a bound that
charges it at rate $\mu$ for the whole of $W$ is bounding a trajectory that does
not occur.

Two readings deserve to be separated. Across the grid, Eq.~\eqref{eq:markov_bound}
is informative essentially only at $\lambda = 50$ under the first-swap window,
and at $\lambda = 25$ for $\Delta e \ge 0.436$. That is a narrow domain --- and
it contains the operating point at which the starvation certificate is stated.
The proposition is not vacuous where it is used; it is vacuous almost everywhere
else, and the submitted version said neither.

\subsection{\texorpdfstring{$W$}{W} is a random variable}\label{sec:prop3_random}

Reviewer~\#1 faults Proposition~9 for treating a random window as a constant.
Proposition~\ref{prop:starvation} is stated at a single $W$, and the campaign
measures a distribution whose spread straddles $\partial\mathcal{D}$. The
statement that bounds the ensemble of runs is therefore

\begin{equation}\label{eq:eq4_integrated}
  \mathbb{P}(S^* \ge \lambda)
  \;\le\;
  \mathbb{E}_W\Bigl[\min\bigl\{1,\, W e^{-\frac{2}{W}(\lambda - s_0 - \mu W)_+^2}\bigr\}\Bigr],
\end{equation}
the expectation taken over the law of $W$ at that magnitude, by conditioning on
$W$ and applying Eq.~\eqref{eq:markov_bound} pathwise.

\begin{remark}[The cap is not cosmetic, and it fixes a floor]\label{rem:eq4_cap}
  A run whose transient exceeds the horizon contributes $T_h = \SixHorizon$ to
  the uncapped integrand, so an uncapped ``bound'' reports $2500$ on a
  probability. Capped, such a run contributes exactly $1$, and
  \eqref{eq:eq4_integrated} is bounded below by the censoring fraction alone:
  $\mathbb{E}_W[\cdot] \ge \mathbb{P}(W \ge T_h)$ at every $\lambda$. On the
  arm integrated over, that fraction is $8.21\%$ ($162$ of $1{,}974$
  complete-case runs; $8.20\%$, $164$ of $2{,}000$, on the whole arm). Any
  claim that \eqref{eq:eq4_integrated} is small at some $\lambda$ is therefore
  bounded away from truth by the censoring before the drift magnitude is even
  consulted.
\end{remark}

Two estimators bracket the integral. The plug-in holds censored runs at their
censoring time and so understates $W$ and the bound; the parametric estimator
integrates the log-normal accelerated-failure-time law already fitted to these
traces with exact right censoring, and so extrapolates past the horizon. Over
the twenty magnitudes:

\smallskip
\centerline{\begin{tabular}{rccc}
  \toprule
  $\lambda$ & deterministic $\bar{W}$ & plug-in & parametric \\
  \midrule
   8 & $1.000$ & $1.000$ & $0.9994$--$1.000$ \\
  15 & $1.000$ & $1.000$ & $0.9902$--$1.000$ \\
  25 & $1.000$ & $1.000$ & $0.9487$--$0.9998$ \\
  50 & $1.000$ & $1.000$ & $0.7526$--$0.9951$ \\
  \bottomrule
\end{tabular}}
\smallskip

\noindent Treating $W$ as random does not repair Eq.~\eqref{eq:markov_bound};
it makes its vacancy unconditional. The deterministic and plug-in readings are
vacant at every magnitude and every threshold. The parametric reading is
formally non-vacant, and its best value over the whole grid is $0.753$: a bound
of $0.75$ on a crossing probability whose measured value is $0.000$ at
$\lambda = 50$. We report that as what it is --- the bound is not wrong, it is
empty --- and we do not present it as a repair.

\subsection{What actually carries the published result}\label{sec:prop3_loadbearing}

Proposition~\ref{prop:certificate} is deterministic. Its chain
$S_{\max}(H) \ge A_{\mathrm{swap}} \ge A_{\mathrm{sig}}(H) - H\delta_P$ uses no
distributional assumption, no rate model and no window length; it is an
inequality between three quantities measured on the same trajectory. Both
Definition~\ref{def:decoupling}(i-bis) and Empirical
Result~\ref{res:tension} route through $A_{\mathrm{swap}}$, and the
synchronised campaign states outright that $W^*$ is the wrong sufficient
statistic: at $\Delta e = 0.3268$, $\kappa$ exceeds its critical value by a
factor of $6$ to $18$ on $100\%$ of runs --- the gate that was supposed to
falsify the $\lambda = 50$ certificate --- and the certificate nonetheless holds,
with $0$ detections in $100$ runs, because the ceiling
$\max_t A_{\mathrm{unrefl}} = 33.5$ is simply below $\lambda = 50$ however long
the transient lasts.

Proposition~\ref{prop:certificate} is therefore the load-bearing statement.
Proposition~\ref{prop:starvation} is a scope statement: it says where a
fluctuation argument can still say something about the crossing probability,
and \S\ref{sec:prop3_domain} says that this is a small set. The two are not in
competition and neither subsumes the other.

\paragraph{The removal branch, argued on both sides.}
Keeping Eq.~\eqref{eq:markov_bound} costs a proposition that is vacant at the
measured operating point under two of its three window estimators, and invites
a referee to ask why a bound that returns $\ge 1$ on $188$ of $240$ measured
configurations is in the paper at all. Removing it costs the only statement
that quantifies the fluctuation margin --- the $\sqrt{W \ln(W/\varepsilon)}$
term, which is what explains why the empirical crossover is a finite-width
stochastic transition rather than a sharp threshold --- and the only bridge to
the post-transient reading of Remark~\ref{rem:sufficient}.

The rule fixed before the domain table was computed retains it on a
non-empty informative domain and withdraws it on an empty one. The domain is
not empty: $32$ of $240$ measured pairs, including the $(\lambda = 50,
W = \tau_{\mathrm{swap}}^{(1/M)})$ pair at which the certificate is stated.
Eq.~\eqref{eq:markov_bound} is therefore \textbf{retained, restricted to
$\mathcal{D}$, and demoted}: it is presented as a scope statement about
fluctuation arguments, with its domain table, and never as the mechanism of the
blind spot. The mechanism is Proposition~\ref{prop:certificate}.

\subsection{Proofs}\label{sec:prop3_proofs}

\begin{proof}[Proof of Theorem~\ref{thm:floor}, and its tightening]
  The proof in \S\ref{sec:requirement} is correct and is not restated. Two
  points are recorded.

  \emph{The chain rule.} $\mathrm{KL}(P_1 \Vert P_0) = \sum_t \mathbb{E}_1
  \bigl[d(\bar{e}_t^{(1)} \Vert \bar{e}_t^{(0)})\bigr]$ holds for the laws of
  the $W$ observed errors in the conditional form, the expectation being over
  the filtration generated by the past. The product form quoted in the proof is
  the special case of independent increments; under a learning classifier the
  errors are not independent, and the bound $u_t \le \Delta_{\max}$ must be read
  as a bound on the conditional excess $\bar{e}_t(\mathcal{F}_{t-1}) - p_0$,
  which is what $\Delta_{\max}$ is defined to be. With that reading every step
  goes through unchanged, and the theorem stands.

  \emph{The relaxation.} The step $d(x \Vert y) \le (x-y)^2/(y(1-y))$ is the
  $\chi^2$ bound, applied $W$ times. It is unnecessary: $u \mapsto
  d(p_0 + u \Vert p_0)$ is convex on $[0, \Delta_{\max}]$ and vanishes at
  $u = 0$, so the chord through the endpoints gives
  $d(p_0+u \Vert p_0) \le (u/\Delta_{\max})\,d_{\max}$ with $d_{\max} :=
  d(p_0 + \Delta_{\max} \Vert p_0)$, whence $\mathrm{KL}(P_1 \Vert P_0) \le
  (d_{\max}/\Delta_{\max})(A + W\delta_P)$ and Eq.~\eqref{eq:floor_chord}.
  Since $d_{\max} \le \Delta_{\max}^2/(p_0(1-p_0))$ is the same $\chi^2$ bound
  applied once, \eqref{eq:floor_chord} is never weaker
  than~\eqref{eq:floor}. At $p_0 = 0.024$, $\Delta_{\max} = 0.3268$ the ratio
  is $6.74$; over the sensitivity band $[0.015, 0.032]$ it runs from $9.00$ to
  $5.72$. The relaxation is loosest exactly where the measured base rate sits,
  which is why it was invisible at the assumed $p_0 = 0.05$.
  \textbf{Status: PROVED} (as stated), and tightened.
\end{proof}

\begin{proof}[Proof of Proposition~\ref{prop:invariance}]
  Under hypothesis~(i), $A = (\Delta e - \delta_P)W$; substituting
  hypothesis~(ii), $W = K(\Delta e)^{-\alpha}$, gives
  $A = K(\Delta e)^{-\alpha}(\Delta e - \delta_P) =
  K(\Delta e)^{1-\alpha}(1 - \delta_P/\Delta e)$, which is
  Eq.~\eqref{eq:invariance}. At $\alpha = 1$ the prefactor is $1$ and
  $A = K(1 - \delta_P/\Delta e) \to K$ as $\Delta e/\delta_P \to \infty$,
  monotonically from below. \textbf{Status: PROVED UNDER STATED CONDITION.}

  The condition is the whole content of the caveat. Hypothesis~(i) is the
  rectangular surrogate, withdrawn on measurement: $A/A_{\mathrm{rect}}$ spans
  $2.70$ to $-14.12$ and changes sign. Hypothesis~(ii) holds at
  $\hat\alpha \approx 0.98$, and Eq.~\eqref{eq:invariance} then predicts a
  spread of $(\Delta e)^{0.02}$, a factor $1.03$ over
  $\Delta e \in [0.10, 0.50]$, against a measured factor of $1.90$ that is not
  monotone. The proposition is thus a true statement about a model the
  repository has withdrawn, and it explains $3\%$ of a $90\%$ effect. What
  survives is Remark~\ref{rem:invariance_measured}'s counterfactual: the same
  ceiling on the branch where adaptation is suppressed spreads by a factor of
  $124$ and is monotone.
\end{proof}

\begin{proof}[Proof of Corollary~\ref{cor:split}]
  Fix $W$, $\varepsilon$, $n_{\mathrm{stat}}$ and let $\alpha \to 0$. The
  $\varepsilon$-margin $\sqrt{(W/2)\ln(1/\varepsilon)}$ is common to the three
  and constant, so only the $\alpha$-price matters.

  \emph{CUSUM.} $\lambda(\alpha)$ is defined by
  $\alpha = W/\mathrm{ARL}_0(\lambda)$ with Siegmund's
  $\mathrm{ARL}_0(\lambda) = (e^{\theta^* \lambda} - \theta^*\lambda - 1)/
  (\theta^*\delta_P)$. Taking logarithms,
  $\theta^*\lambda = \ln(1/\alpha) + \ln(W\theta^*\delta_P) + \ln(1 +
  O(\theta^*\lambda e^{-\theta^*\lambda}))$, so
  $\lambda(\alpha) = \theta^{*-1}\ln(1/\alpha)\,(1+o(1))$: linear in
  $\ln(1/\alpha)$, with no $W$ in the leading term.

  \emph{ADWIN and KSWIN.} By Eqs.~\eqref{eq:Radwin}--\eqref{eq:Rkswin} the
  $\alpha$-prices are $\sqrt{(W/2)\ln(4W/\alpha)}$ and
  $\sqrt{n_{\mathrm{stat}}\ln(2/\alpha)}$, both
  $\Theta(\sqrt{\ln(1/\alpha)})$ at fixed $W$ and $n_{\mathrm{stat}}$.

  A linear function of $\ln(1/\alpha)$ exceeds any square-root function of it
  beyond a finite crossing point, so for every fixed budget $A$ there is
  $\alpha^\dagger > 0$ below which $R_{\mathrm{CUSUM}} > A \ge
  R_{\mathrm{ADWIN}} \wedge R_{\mathrm{KSWIN}}$ fails for the cumulative monitor
  alone. The $W$-scalings are read directly off the same three expressions:
  $\lambda(\alpha)$ carries no $W$; $\sqrt{(W/2)\ln(4W/\alpha)} \to 0$ as
  $W \to 0$ because $W\ln W \to 0$; and $\sqrt{n_{\mathrm{stat}}\ln(2/\alpha)}$
  carries no $W$ either. \textbf{Status: PROVED.}

  The last of those three corrects the submitted statement, which placed the
  $\alpha$-price of $R_{\mathrm{KSWIN}}$ ``inside a square root of $W$''.
  Eq.~\eqref{eq:Rkswin} contains no such $W$, and the correction matters: KSWIN
  escapes not because its requirement vanishes with the transient but because
  its $\alpha$-price is set by the statistic size, which the practitioner
  chooses, instead of by a threshold that must grow like $\ln(1/\alpha)$.
  Remark~\ref{rem:split_measured} evaluates all three at a common $\alpha$ and
  exhibits the crossing at $\ln(1/\alpha) \approx 13$.
\end{proof}

\subsection{EDDM}\label{sec:prop3_eddm}

No $R_{\mathrm{EDDM}}$ is delivered, and the requirement is withdrawn from the
family of \S\ref{sec:requirement}. The reasons are stated so that the
withdrawal is not read as an oversight.

EDDM compares $p'_i + 2s'_i$, the cumulative mean and standard deviation of
inter-error distances since the last reset, against the running maximum of that
same path. Modelling the distances as geometric at rate $p_0$ before the change
and $p_1$ after, and writing $f$ for the post-change share of errors since the
reset, $d := 1/p$ and $\sigma^2 := (1-p)/p^2$,
\begin{equation}\label{eq:eddm_level}
  \mathrm{level}(f) =
  \frac{(1-f)d_0 + f d_1 + 2\sqrt{(1-f)\sigma_0^2 + f\sigma_1^2
        + f(1-f)(d_0-d_1)^2}}{d_0 + 2\sigma_0},
\end{equation}
decreasing from $\mathrm{level}(0) = 1$, so detection requires
$f \ge f^\dagger := \inf\{f : \mathrm{level}(f) < \beta\}$ and hence
$W_{\mathrm{EDDM}} = (n_0/p_1)\,f^\dagger/(1-f^\dagger)$ post-change steps,
with $n_0$ the pre-change error count since the last reset. The between-group
term $f(1-f)(d_0-d_1)^2$, which the normal approximation on geometric distances
omits, grows with the contrast at the same time as the mean falls, and the two
largely cancel: $f^\dagger$ is nearly invariant in the drift magnitude, $0.235$
to $0.240$ over $p_1 \in [0.10, 0.70]$, so $W_{\mathrm{EDDM}} \approx
0.31\, n_0/p_1$. It therefore shortens with the drift magnitude much as a
first-passage time does. What distinguishes it is the factor $n_0$: a threshold
test requires an \emph{absolute} amount of evidence, whereas EDDM requires a
fixed \emph{share} of its own history, so its requirement grows without bound as
the stationary phase lengthens. No first-passage bound carries such a term.

This closed form does not reproduce the measured collapse, and the rule fixed in
advance makes withdrawal the outcome in that case. At the ProteuS pre-change
history it returns $168$--$281$ steps against the normal approximation's $155$:
the same order of magnitude, not the ``no detection'' that $0$ alarms in
$1{,}080$ runs requires. Worse, it cannot be evaluated where the collapse is
measured: the ProteuS artifacts record $F_1$, ADD and seed and no error stream,
so neither $p_0$ nor $p_1$ nor the transient length is available there.

The one stream on which the inputs do exist cannot arbitrate either. On the
synchronised traces, the detector has not completed its $30$-error warm start at
the change point in $95\%$ of runs --- $1{,}000$ pre-change steps at
$p_0 = 0.024$ furnish about $24$ errors --- so every alarm is post-warm-start by
construction; on the \texttt{static} arm, which does warm up, it alarms
\emph{before} the change point in $60\%$ of runs. The decisive control is that
the three adaptive arms return the same alarm rate to within $0.08$, although
they differ by exactly the adaptation under study, and the \texttt{frozen} arm
never adapts at all. A statistic that cannot tell those three apart is not
reporting on drift.

The \emph{empirical} EDDM result is untouched by this: $0$ detections in
$1{,}080$ runs against $882$ for the same detector on a non-adaptive learner,
sign test $p \approx 1.9\times10^{-9}$, is a measurement and stands on its own.
What is withdrawn is the claim that a requirement $R_{\mathrm{EDDM}}$ of the
form of Eqs.~\eqref{eq:Rcusum}--\eqref{eq:Rkswin} governs it.

\subsection{Proof status}\label{sec:prop3_status}

\centerline{\begin{tabular}{llp{0.42\linewidth}}
  \toprule
  Statement & Status & Condition \\
  \midrule
  Thm.~\ref{thm:floor} & PROVED & chain rule read conditionally; tightened by
    Eq.~\eqref{eq:floor_chord} \\
  Prop.~\ref{prop:invariance} & PROVED UNDER STATED CONDITION &
    rectangular transient and $\alpha = 1$; both measured false \\
  Cor.~\ref{cor:split} & PROVED & at fixed $W$, $\varepsilon$,
    $n_{\mathrm{stat}}$, as $\alpha \to 0$ \\
  Prop.~\ref{prop:order} & PROVED & drift not resolved before the first
    replacement; holds on $1{,}835/1{,}836$ estimable runs \\
  $R_{\mathrm{EDDM}}$ & NOT PRODUCED & withdrawn; the empirical result stands \\
  \bottomrule
\end{tabular}}
